\section*{Benchmark Construction}
NeuroscienceBench comprises 91 evaluations drawn from 43 peer-reviewed publications and their public datasets, which we refer to as problems (Figure~\ref{fig:anatomy}). A problem yields several evaluations when its dataset supports more than one independent decision. The evaluations cover the field broadly along all three taxonomy axes.

The functional-system axis records which nervous-system function is studied, such as motor, vision, olfaction \& taste, somatosensation, language \& social and others. The 33 cross-system evaluations address a tissue compartment or whole-person readout rather than a single pathway, most often cortical and forebrain neurons, vasculature, barriers and CSF, and glia and neuroimmune. The process-class axis separates basic-science questions about how the nervous system normally works from questions about how it fails.  Normal function is subdivided into eight subcategories, such as brain state and sleep, cognition, and motor control, and disease, injury and aging into eleven process classes, such as neurodegeneration, infection, and aging. The biological-level axis records the scale at which the evaluation's decision is made, from molecular through cellular and circuit to tissue, organ, and organism behavior.

Each evaluation also carries a workflow-stage tag for one of six phases of the research cycle: hypothesis and target, model-system selection, experimental design, data processing, analysis and inference, and validation and translation.  The underlying data include next-generation sequencing, extracellular electrophysiology, histology and cellular imaging, sleep-state scoring, BCI decoding, EEG and ECoG, behavioral scoring and pose tracking, patch clamp and multi-electrode arrays, and clinical and biofluid cohorts.

Evaluations were built with a multi-stage authoring pipeline. Before inclusion, each evaluation passed automated checks for structure, shortcuts (tasks answerable without analyzing the data), guessability and near-duplication. Grading is deterministic. Numeric answers must fall within a tolerance or admissible range, categorical answers must match exactly, and sets are scored by overlap.
